\subsection{Barrido de Eventos (Máximo Solapamiento de Intervalos)}
\label{alg:5-5}

\graybold{Preguntas guía:}

\begin{itemize}
\item ¿Tengo intervalos (entrada y salida, inicio y fin, reserva y devolución) y me piden cuántos coinciden como máximo en algún instante?
\item ¿Me sirve olvidar qué intervalo es cada uno y quedarme solo con los EVENTOS ordenados en el tiempo?
\item ¿El estado cambia únicamente en los extremos de los intervalos, de modo que basta revisar esos instantes?
\item ¿Los tiempos son enormes (\ten{9}) pero la cantidad de eventos es manejable, descartando recorrer la línea de tiempo entera?
\end{itemize}

\graybold{Utilidad:} El barrido de eventos convierte un problema sobre intervalos en un problema sobre una lista ordenada de puntos: cada intervalo se descompone en un evento de apertura y uno de cierre, se ordenan todos por tiempo y se recorre llevando un contador. Resuelve máximo solapamiento (salas necesarias, clientes simultáneos, canales de un horario), unión de longitudes cubiertas, y es la base de los barridos en geometría. La clave es que el contador solo cambia en los extremos, así que los \ten{9} instantes posibles se reducen a $2n$ puntos relevantes.

\graybold{Complejidad:} \bigO{n \log n} dominado por el ordenamiento de los $2n$ eventos, más \bigO{n} para el recorrido.

\graybold{Fundamento teórico:} Sea $f(t)$ la cantidad de intervalos que contienen al instante $t$. Esa función solo cambia en los extremos: sube en uno en cada llegada y baja en uno en cada salida, y entre dos eventos consecutivos es constante. Por lo tanto su máximo se alcanza inmediatamente después de alguna llegada, y basta recorrer los eventos en orden cronológico acumulando $+1$ y $-1$ y registrando el máximo del acumulado. El enunciado garantiza que todos los tiempos son DISTINTOS, lo que evita la única ambigüedad del método: cuando una salida y una llegada coinciden en el mismo instante, hay que decidir si el cliente que se va todavía cuenta, y según el problema conviene procesar primero una o la otra (si los intervalos son cerrados y tocarse cuenta como coincidir, la llegada va primero; si son semiabiertos, la salida). Aquí, en lugar de ordenar pares, cada evento se codifica en un solo entero, $2t+1$ para una llegada y $2t$ para una salida: al duplicar el tiempo, el bit bajo queda libre para marcar el tipo sin alterar el orden temporal, ordenar una lista de enteros es más rápido que ordenar tuplas, y la desigualdad $2t < 2t+1$ deja a las salidas antes que las llegadas del mismo instante, que es el criterio conservador. El bit se lee después con un \texttt{and} para saber si sumar o restar.~\cite{shamos1976}

\graybold{Ejercicio aplicado}

(CSES 1619 — Restaurant Customers) Dadas las horas de llegada y salida de $n$ clientes de un restaurante, hallar la máxima cantidad de clientes presentes simultáneamente.

\graybold{I/O:} $n \le 2\cdot10^5$ con dos enteros por cliente $\Rightarrow$ lectura total + tokenización (patrón 2), separando llegadas y salidas con slices de paso 2 y codificándolas directamente al construir la lista, sin pasar por una estructura intermedia. Verificado contra el caso de ejemplo y contrastado con una fuerza bruta que, para cada instante relevante, cuenta cuántos intervalos lo contienen, en 500 casos aleatorios con tiempos todos distintos como garantiza el enunciado, sin discrepancias. Con $n = 2\cdot10^5$ tarda entre \aprox{}0.08s (intervalos todos anidados o todos disjuntos, donde el ordenamiento encuentra datos casi ordenados) y \aprox{}0.20s (tiempos aleatorios); el límite es de 1 segundo.

\graybold{Código solución}

\begin{lstlisting}[language=Python]
import sys

def solve():
    data = sys.stdin.buffer.read().split()
    n = int(data[0])
    # cada evento se codifica en un entero: llegada = 2t+1, salida = 2t.
    # duplicar el tiempo deja libre el bit bajo para marcar el tipo sin
    # alterar el orden cronologico, y ordenar enteros es mas rapido que tuplas
    ev = [2*int(v) + 1 for v in data[1:1 + 2*n:2]]     # llegadas
    ev += [2*int(v) for v in data[2:2 + 2*n:2]]        # salidas
    ev.sort()
    cur = 0
    mejor = 0
    for v in ev:
        if v & 1:                # bit bajo encendido: es una llegada
            cur += 1
            if cur > mejor: mejor = cur
        else:
            cur -= 1
    print(mejor)

solve()
\end{lstlisting}
